
\subsubsection{3.1 Dataset: h-e2 Survival Panel}

The analysis is based on the \textbf{h-e2 panel}, a validated survival dataset constructed from Papers With Code \cite{paperwithcode2025}. The panel covers 87 parent tasks from the Koch et al. \cite{koch2021reduced} 87-task taxonomy, spanning 2015--2023. Each row corresponds to a benchmark that achieved plurality status (highest SOTA score for its task) at some point in the observation window.

\textbf{Panel statistics:} 345 total rows; 87 tasks; temporal window 2015--2023; validated across 10 prior pipeline attempts (14/14 unit tests pass).

\textbf{Survival variables:} \texttt{duration} (years the benchmark held plurality status); \texttt{event} = 1 if displaced, 0 if right-censored.

\textbf{Complete-case analysis:} 87 rows with missing predictor values were excluded, yielding 258 complete-case rows. Of these, 176 are displacement events (event = 1) and 82 are right-censored (event = 0); event rate = 68.2\%.

\textbf{Baseline covariates:} \texttt{task\textbackslash \_age} (years since task appeared in PWC); \texttt{log\textbackslash \_publication\textbackslash \_volume} (log total paper count for task). Note: \texttt{benchmark\textbackslash \_introduction\textbackslash \_year} was excluded due to perfect collinearity with \texttt{task\textbackslash \_age} (|r| = 1.000), detected automatically by Gate G4.

\subsubsection{3.2 Diversity Predictor Construction}

Diversity predictors were constructed from \textbf{pwc-archive/evaluation-tables} \cite{paperwithcode2025} (HuggingFace, CC-BY-SA-4.0, July 2025 snapshot; 59,860 rows after PyArrow flattening from the nested 326k-row structure). The key field is \texttt{paper\textbackslash \_url}, identifying distinct paper submissions per benchmark.

\textbf{Primary predictor:} \texttt{log\textbackslash \_unique\textbackslash \_paper\textbackslash \_count\textbackslash \_at\textbackslash \_intro\textbackslash \_z} = log1p(unique \texttt{paper\textbackslash \_url} count per benchmark through introduction year), z-standardized. This measures breadth of community participation in paper submissions at introduction.

\textbf{Secondary predictor:} \texttt{paper\textbackslash \_diversity\textbackslash \_ratio\textbackslash \_at\textbackslash \_intro\textbackslash \_z} = (unique count / total submission rows), z-standardized. Normalizes by total volume, removing time-accumulation effects.

\textbf{Join:} Task paths were matched via rapidfuzz token\_sort\_ratio (threshold = 85), achieving 86.2\% coverage (75 of 87 benchmarks matched). Diversity statistics were computed for 67 of 87 benchmarks. Predictor collinearity: Pearson r = $-0.324$ between the two diversity predictors --- independent constructs.

\subsubsection{3.3 5-Gate FAIL FAST Pre-Validation Protocol}

A key methodological contribution is the \textbf{FAIL FAST protocol} (H-E1), which validates data quality \textit{before} any Cox regression. Gates are evaluated in sequence; failure at any gate routes to hypothesis revision rather than analysis.

\begin{table*}[htbp]
\centering
\resizebox{\dimexpr 2\columnwidth+\columnsep\relax}{!}{%
\begin{tabular}{lllll}
\toprule
Gate & Name & Criterion & Measured Value & Status \\
\midrule
G0 & Join Coverage & $\geq$ 80\% benchmarks with non-null paper\_url & 0.862 & PASS \\
G1 & Primary Time-Independence & partial $r^{2}(log$\_count\_z ~ temporal) > 0.01 & 0.605 & PASS \\
G2 & Secondary Time-Independence & partial $r^{2}(diversity$\_ratio\_z ~ temporal) > 0.01 & 0.975 & PASS \\
G3 & Variance & std(diversity\_ratio) > 0.10 & 0.246 & PASS \\
G4 & Multicollinearity & max VIF < 10 & 2.14 & PASS \\
\bottomrule
\end{tabular}
}
\end{table*}

\textbf{Partial $R^2$ computation:} The predictor is regressed on \texttt{[task\textbackslash \_age, benchmark\textbackslash \_introduction\textbackslash \_year]} via OLS; $R^2$ of the residuals is computed. Partial $R^2$ substantially exceeding 0.01 indicates the predictor retains genuine cross-benchmark information beyond temporal trend.

\begin{figure}[htbp]
\centering
\includegraphics[width=\columnwidth]{figures/gate_metrics.png}
\caption{FAIL FAST Gate Results}
\label{fig:gate_metrics}
\end{figure}

\textit{Figure 1: 5-gate FAIL FAST results. All five gates pass by wide margins.}

\begin{figure}[htbp]
\centering
\includegraphics[width=\columnwidth]{figures/partial_r2.png}
\caption{Partial $R^2$ Time-Independence}
\label{fig:partial_r2}
\end{figure}

\textit{Figure 2: Partial $R^2$ values confirming time-independence of diversity predictors (G1 = 0.605, G2 = 0.975), both far exceeding the 0.01 threshold.}

\subsubsection{3.4 Cox Proportional Hazards Model}

\textbf{CoxPHFitter(penalizer=0.1)} from lifelines \cite{davidson2019lifelines} is used throughout. The L2 penalty (penalizer = 0.1) was selected for convergence stability; it was not optimized for effect estimation (see Limitation L5).

\textbf{M0 (baseline):} h(t) = $h_{0}(t) \cdot$ exp($\beta _{1}\cdot task$\_age + $\beta _{2}\cdot log$\_publication\_volume)

\textbf{M1 (diversity predictor):} h(t) = $h_{0}(t) \cdot$ exp($\beta _{1}\cdot task$\_age + $\beta _{2}\cdot log$\_publication\_volume + $\beta _{3}\cdot log$\_unique\_paper\_count\_at\_intro\_z)

\textbf{Likelihood ratio test (LRT):} LRT stat = $-2\cdot (\ell (M0) - \ell (M1))$, df = 1; p-value from chi-squared distribution (scipy chi2.sf; negative-statistic guard applied).

\textbf{Pre-specified success criteria:} P1 (primary): LRT p < 0.05 AND |$HR-1| \geq$ 0.10.

\textbf{Pre-specified direction protocol:}
\begin{itemize}
\item HR < 1, p < 0.05 $\rightarrow$ H1 lock-in supported
\item HR > 1, p < 0.05 $\rightarrow$ H2 saturation supported
\item p > 0.05 (with G0--G4 passing) $\rightarrow$ H0 null
\end{itemize}

